%!TeX root=MemoriaTFG.tex

\chapter{Ablación: modelo fundacional inglés (DermLIP) con
traducción ES--EN}
\label{cap:anexoi}

Este anexo documenta una ablación independiente que cuantifica el
\emph{trade-off} entre la rama contrastiva de Dermapixel R0
defendida en la
sección~\ref{sec:spanderm-clip} y un \emph{wrapper} alternativo
basado en un modelo fundacional dermatológico inglés (DermLIP v2)
con traducción de etiquetas ES--EN, sobre el mismo conjunto de test
\emph{holdout} de DermapixelAI~1.0.

\section{Motivación}
\label{sec:anexoi-motivacion}

Una pregunta legítima planteada durante la revisión paper-grade
de la rama contrastiva de Dermapixel R0 es si el esfuerzo experimental
documentado en la sección~\ref{sec:spanderm-clip} (cinco
iteraciones de entrenamiento cross-modal castellano) se justifica
frente a una alternativa pragmática: emplear un modelo fundacional
dermatológico inglés ya alineado imagen-texto, sin entrenar ningún
modelo adicional, sobre las etiquetas castellanas traducidas al
inglés. Esta ablación cuantifica el \emph{trade-off} entre ambas
estrategias sobre el mismo conjunto de test \emph{holdout} de
DermapixelAI~1.0.

\section{Protocolo experimental}
\label{sec:anexoi-protocolo}

El modelo fundacional inglés evaluado es DermLIP v2
(\texttt{PanDerm-base-w-PubMed-256}), un sistema CLIP-style con
codificador visual PanDerm-Base y codificador textual PubMedBERT,
preentrenado contrastivamente sobre el dataset
Derm1M~\cite{derm1m} ($\approx 403\,000$ pares imagen-texto).
Ambas torres permanecen congeladas; no se entrena ningún
parámetro. Las $54$ clases L3 representadas en el test se
tradujeron al inglés clínico mediante un glosario estándar
(Bolognia/Rook) con verificación cruzada de unicidad, y por cada
clase se generó un \emph{rich prompt} clínico en inglés ($\sim 37$
palabras, longitud equiparable a los \emph{prompts} castellanos
de Dermapixel R0). La inferencia replica el protocolo híbrido de
la sección~\ref{sec:spanderm-clip} (prototipos visuales sobre
\emph{train} más texto rico, régimen \emph{H6} $w_v = 0{,}3$,
$w_t = 0{,}7$) sobre las mismas $101$ imágenes de test y $54$
clases.

Adicionalmente, se midió el \emph{eje idioma puro} sobre el propio
\emph{checkpoint} v4 de Dermapixel R0: el mismo modelo evaluado con
los \emph{prompts} castellanos originales frente a los mismos
\emph{prompts} traducidos al inglés, aislando el efecto del idioma
de consulta sobre una alineación cross-modal fija.

\section{Resultados}
\label{sec:anexoi-resultados}

La tabla~\ref{tab:ablation_panderm_en} reúne las seis
configuraciones evaluadas sobre el mismo conjunto de test.

\begin{table}[H]
\centering
\footnotesize
\setlength{\tabcolsep}{4pt}
\renewcommand{\arraystretch}{0.95}
\caption{Comparativa cuantitativa de Dermapixel R0 v4 (\emph{headline})
frente a la ablación con modelo fundacional inglés DermLIP v2 más
traducción ES--EN. Test \emph{holdout} DermapixelAI~1.0,
$N = 101$, $54$ clases L3 representadas; régimen \emph{H6 hybrid}
$v30\_r70$ salvo donde se indica.}
\label{tab:ablation_panderm_en}
\begin{tabular}{llrrr}
\hline
Sistema & Régimen & acc@1 L3 & BAcc L3 & acc@5 L3 \\
\hline
Dermapixel R0 v4    & H6 hybrid (castellano)            & $0{,}2475$          & $0{,}2685$ & $0{,}4455$ \\
Dermapixel R0 v5    & H6 hybrid (castellano)            & $0{,}2079$          & $0{,}2185$ & $0{,}4158$ \\
Dermapixel R0 v4    & H6 hybrid (inglés)              & $0{,}2277$          & $0{,}2556$ & $0{,}4554$ \\
DermLIP v2          & H6 hybrid (inglés)              & $0{,}2871$          & $0{,}3667$ & $0{,}6436$ \\
\textbf{DermLIP v2} & \textbf{templates+rich (inglés)} & $\mathbf{0{,}3861}$ & $0{,}4099$ & $0{,}6931$ \\
DermLIP v2          & visual-only                     & $0{,}2178$          & $0{,}2278$ & $0{,}4851$ \\
\hline
\end{tabular}
\end{table}

\section{Interpretación}
\label{sec:anexoi-interpretacion}

DermLIP v2 inglés supera a Dermapixel R0 v4 en el régimen
comparable \emph{H6} ($0{,}2871$ frente a $0{,}2475$,
$+3{,}96$\,pp) y lo supera con mayor margen en su mejor régimen
de plantillas más texto rico ($0{,}3861$, $+13{,}86$\,pp),
confirmando la hipótesis de partida de que el cuello de botella de
Dermapixel R0 residía en la \emph{escala} de la alineación
cross-modal y no en la formulación textual. Un modelo fundacional
preentrenado sobre $\approx 403\,000$ pares supera a un CLIP
castellano dedicado entrenado sobre $854$ pares, incluso pagando el
ruido de la traducción de etiquetas.

El eje idioma matiza esta lectura: sobre el mismo \emph{checkpoint}
v4, la consulta en castellano supera a la consulta en inglés
($0{,}2475$ frente a $0{,}2277$ en \emph{H6}, $+1{,}98$\,pp;
y $0{,}1287$ frente a $0{,}0891$ en \emph{rich-only},
$+3{,}96$\,pp), mientras que el régimen visual-only permanece
idéntico ($0{,}1683$). Dermapixel R0 sí captura un componente
específico del registro clínico castellano, si bien de segundo orden
frente a la ventaja de escala del modelo fundacional inglés.

El \emph{trade-off} entre los dos enfoques es, por tanto,
explícito:

\begin{itemize}
  \item \textbf{DermLIP inglés más traducción}: rendimiento bruto
        superior, cero entrenamiento, pero dependencia de
        traducción de etiquetas y de un dataset de preentrenamiento
        externo (Derm1M~\cite{derm1m}), con la consiguiente
        pérdida potencial de matices clínicos castellanos no
        traducibles uno a uno.
  \item \textbf{Dermapixel R0 (rama contrastiva)}: entrenamiento dedicado sobre $854$
        pares reales castellanos, preservación de la
        \emph{fingerprint} clínica del archivo de la
        Dra.~R.~Taberner, inferencia sin traducción en
        \emph{runtime} y operación \emph{on-prem} compatible con la
        ruta de extensión al banco hospitalario HUSLL
        (sección~\ref{sec:linea-active-learning}).
\end{itemize}

La conclusión es que la rama contrastiva de Dermapixel R0 no se
justifica como sistema de
mayor rendimiento bruto en \emph{zero-shot} L3, sino por sus
propiedades operativas de privacidad, especificidad lingüística y
autonomía de despliegue. La superioridad de DermLIP refuerza,
además, el diseño de la línea de continuación documentada en la
sección~\ref{sec:linea-active-learning}: la palanca pendiente
para Dermapixel R0 es la anotación experta dirigida
\emph{in-domain} sobre las clases L3 con peor cobertura, no la
escala genérica.

\section{Limitaciones de la ablación}
\label{sec:anexoi-limitaciones}

La comparación no está equiparada en datos: DermLIP v2 se
preentrena sobre $\approx 403\,000$ pares frente a los $854$ pares
castellanos de Dermapixel R0, y los \emph{backbones} visuales
difieren (PanDerm-Base contrastivo frente a PanDerm-Large
congelado con LoRA). El resultado debe leerse como ``modelo
fundacional inglés grande más traducción'' frente a ``CLIP
castellano dedicado pequeño'', no como un contraste de igual
presupuesto de datos.

La carga de DermLIP v2 requiere el \emph{fork} de
\texttt{open\_clip} del proyecto Derm1M: la implementación estándar
instancia un \emph{Vision Transformer} genérico e ignora
silenciosamente la torre visual PanDerm-Base ($152$ claves
ausentes, codificador visual aleatorio), un modo de fallo
verificado y evitado en esta ablación empleando el mismo
\emph{loader} que el módulo de producción
(Anexo~\ref{cap:anexoh}).

La verificación de solapamiento con Derm1M se realizó a nivel de
fuente: las fuentes constituyentes de Derm1M (vídeo, foros,
literatura biomédica y libros de texto en inglés y chino) no
incluyen ningún blog dermatológico español, y no se hallaron
coincidencias para \emph{dermapixel} ni \emph{taberner} en las
$417\,257$ filas del \emph{manifest}. El archivo Dermapixel es
privado y no figura como fuente, lo que sustenta razonablemente la
asunción de ausencia de \emph{leakage}; no se ejecutó una
comparación perceptual exhaustiva imagen a imagen. La traducción
ES--EN introduce, finalmente, un ruido no cuantificado en las
etiquetas.

\section{Detalles de reproducibilidad de la rama contrastiva}
\label{sec:anexoi-repro-contrastiva}

Este apartado recoge los detalles de procedencia y reproducibilidad
de la rama contrastiva de Dermapixel R0
(sección~\ref{sec:spanderm-clip}), separados del cuerpo por no
afectar a las conclusiones científicas.

\paragraph{Generación de \emph{captions} sintéticas.}
Las ocho perspectivas clínicas por imagen se generan mediante un
modelo de lenguaje grande a partir del texto original de cada caso
del blog y de sus metadatos. El total alcanza $7\,680$
\emph{captions} sobre $960$ imágenes (fase v1--v4); $102$ imágenes
fallan la generación automática y conservan únicamente el texto
base. El coste real de generación es de $32{,}54$ USD ($0{,}030$ USD
por imagen completa) y la longitud media es de $87{,}5$ palabras por
\emph{caption}. Combinadas con las nueve variantes textuales base
(plantillas L1/L2/L3, jerárquicas \texttt{L1L2L3}, título de blog,
tres \emph{chunks} de \texttt{case\_text} y \texttt{diagnosis\_raw}),
el dataset de entrenamiento alcanza del orden de $22\,000$ pares
(imagen, texto), muestreados por hash determinista
\texttt{(seed, epoch, image\_id)}.

\paragraph{Tiempos de entrenamiento.}
Sobre NVIDIA DGX Spark GB10: $\sim 18$\,min (v1), $\sim 3$\,min
(v2), $\sim 5$\,min (v3), $\sim 9$\,min (v4) y $\sim 54$\,min (v5).

\paragraph{Lección metodológica: longitud de las \emph{captions}.}
La generación de \emph{captions} sintéticas del dataset externo de
v5 (coste $42{,}33$ USD) reveló la importancia de calibrar la
longitud objetivo respecto al umbral del verificador automático
posterior. Un primer intento con rango $30$--$40$ palabras produjo
una tasa de rechazo del $54\%$ ---las \emph{captions} se agrupaban
en el suelo del rango y el verificador descartaba las más cortas---;
recalibrar el objetivo a $40$--$55$ palabras redujo el rechazo a
cero. La regla operativa derivada es dimensionar la longitud
objetivo con margen explícito sobre el umbral mínimo aceptable, no
sobre la media deseada.

\section{Material complementario de la rama contrastiva}
\label{sec:anexoi-material-contrastivo}

Este apartado recoge las tablas y figuras de detalle de la
sección~\ref{sec:spanderm-clip}, desplazadas del cuerpo por
constituir diagnóstico secundario sin alterar las conclusiones
científicas allí defendidas.

\subsection{Cabeza supervisada L2: \emph{best-checkpoint} frente a
últimas épocas}
\label{sec:anexoi-l2-checkpoint}

La cabeza supervisada L2 de Dermapixel R0
(sección~\ref{sec:dermapixel-spanderm-v0}) se reporta sobre el test L2
($N_{\mathrm{test}} = 36$) con dos estimaciones: el \emph{best val
checkpoint} (selección por máxima BAcc de validación) y la media de las
últimas cinco épocas sin selección, esta última como comportamiento
asintótico sin la varianza del criterio de selección
(tabla~\ref{tab:dermapixel-spanderm-l2}).

\begin{table}[H]
\centering
\footnotesize
\setlength{\tabcolsep}{4pt}
\renewcommand{\arraystretch}{0.95}
\caption{Métricas de Dermapixel R0 (cabeza supervisada L2) sobre el
conjunto de test L2 de
DermapixelAI~1.0 ($N_{\mathrm{test}} = 36$, $38$ clases con
queratinización consolidada). Cifras: media $\pm$ desviación
estándar sobre tres semillas ($\{42, 43, 44\}$). La columna
\emph{best val} aplica selección por máxima BAcc de validación;
la columna \emph{últimas $5$ ép.} promedia las cinco épocas
finales sin selección. Mejor por fila en negrita.}
\label{tab:dermapixel-spanderm-l2}
\begin{tabular}{lrr}
\hline
Métrica   & \emph{Best val \emph{checkpoint}} & Media últimas $5$ ép. \\
\hline
Acc@1     & $\mathbf{0{,}250 \pm 0{,}000}$ & $0{,}209 \pm 0{,}005$ \\
Acc@3     & $\mathbf{0{,}500 \pm 0{,}023}$ & $0{,}457 \pm 0{,}011$ \\
BAcc      & $\mathbf{0{,}363 \pm 0{,}007}$ & $0{,}309 \pm 0{,}006$ \\
AUROC     & $\mathbf{0{,}855 \pm 0{,}006}$ & $0{,}837 \pm 0{,}008$ \\
W-F1      & $\mathbf{0{,}218 \pm 0{,}013}$ & --- \\
Kappa     & $\mathbf{0{,}220 \pm 0{,}001}$ & --- \\
\hline
\end{tabular}
\end{table}

La estabilidad entre semillas es alta (desviación $\leq 0{,}013$,
$\leq 0{,}008$ sobre AUROC y BAcc). La diferencia entre el \emph{best
val checkpoint} y la media de las últimas cinco épocas ($+5{,}4$\,pp
BAcc, $+1{,}8$\,pp AUROC, $+4{,}1$\,pp Acc@1) refleja un sobreajuste
suave análogo al del FT parcial denso
(sección~\ref{sec:dermapixel-ft-l1}), atenuado por la restricción al
$0{,}17\%$ de parámetros entrenables de LoRA.

\subsection{Barrido de épocas de la iteración v3}
\label{sec:anexoi-epoch-sweep}

\begin{figure}[H]
\centering
\includegraphics[width=0.88\textwidth]{fig_spanderm_epoch_sweep.png}
\caption{Barrido de épocas de la iteración contrastiva v3 de
Dermapixel R0 sobre el
conjunto de test de $101$ imágenes. \emph{Accuracy}@1 L1, L2 y L3
y \texttt{val i2t\_R@5} frente al paso de entrenamiento. L3
oscila en el intervalo $[0{,}04,\,0{,}12]$ a lo largo de las diez
épocas sin aproximarse al objetivo $\geq 0{,}25$, mientras
que \texttt{val i2t\_R@5} alcanza su máximo en la época $2$ y
desciende con el entrenamiento conjunto adicional del codificador
visual.}
\label{fig:spanderm-epoch-sweep}
\end{figure}

\subsection{Techo de \emph{Linear Probing} y diagnóstico mecanístico}
\label{sec:anexoi-lp-ceiling}

Para entender por qué ningún cambio textual mueve el \emph{zero-shot}
L3 se mide el techo \emph{Linear Probing} sobre las \emph{features}
$1\,024$-D del encoder visual, con script idéntico para v2 (frozen) y
v3 (LP en épocas $2$ y $6$). La tabla~\ref{tab:spanderm-lp-ceiling}
reúne los tres puntos de la comparación.

\begin{table}[H]
\centering
\footnotesize
\setlength{\tabcolsep}{4pt}
\renewcommand{\arraystretch}{0.95}
\caption{Techo \emph{Linear Probing} supervisado sobre
\emph{features} $1\,024$-D del image encoder. Comparativa
\emph{frozen} (v2 e$7$) frente a image-LoRA tras entrenamiento
adicional (v3 época $6$). Métrica: \emph{accuracy}@1. Las
cifras de v3 corresponden a la evaluación intermedia conservada
del Sprint97; la reejecución posterior arroja desviaciones
$< 0{,}01$ absolutas, consistentes con ruido de muestreo
($N = 101$) y sin alterar la lectura ($+0{,}01 \approx$ ruido).}
\label{tab:spanderm-lp-ceiling}
\begin{tabular}{lrrrr}
\hline
Nivel              & LP v2-e$7$ frozen & LP v3 best.pt (e$2$) & LP v3 e$6$ ($+$image train) & $\Delta$ frozen$\to$trained \\
\hline
L1 ($4$ vías)       & $0{,}584$         & $0{,}584$            & $0{,}594$                   & $+0{,}01$                   \\
L2 ($38$ vías)      & $0{,}214$         & $0{,}214$            & $0{,}225$                   & $+0{,}01$                   \\
L3 ($225$ vías)     & $0{,}079$         & $0{,}079$            & $0{,}089$                   & $+0{,}01$                   \\
\hline
\end{tabular}
\end{table}

Dos observaciones estructuran el diagnóstico. Primero, el
\emph{zero-shot} L3 de v1/v2/v3 ($\sim 0{,}13$) supera ya el techo LP
supervisado ($0{,}079$--$0{,}089$) sobre el mismo encoder frozen: las
\emph{features} $1\,024$-D del PanDerm-Large congelado no separan
linealmente los diagnósticos castellanos finos, y ningún ajuste textual
puede extraer información que las \emph{features} visuales no contienen.
Segundo, la LoRA visual de v3 (últimos cuatro bloques, $1{,}05$\,M
parámetros) eleva el techo LP sólo $+0{,}01$ por nivel, dentro del
ruido: capacidad insuficiente para reesculpir las \emph{features} en el
régimen \emph{fine-grained} castellano (veredicto \emph{``image LoRA
did not help''}: no destruye el preentrenamiento ---el LP nunca cae---
pero tampoco lo mejora).

El mecanismo por el que el \emph{zero-shot} de v3 cae sin que el LP
suba es que el entrenamiento conjunto desplaza los embeddings visuales
hacia el espacio de captions ricas y los aleja del de plantillas de
nombre de clase: L1 \emph{zero-shot} \emph{accuracy}@1 cae
$0{,}426$ (e$2$) $\to 0{,}287$ (e$7$) mientras el \emph{train loss}
sigue bajando a $0{,}29$, señal de sobreajuste \emph{rich-caption}. La
capacidad arquitectónica (v4) es la palanca efectiva; la escala externa
genérica (v5) no. En corpora clínicos pequeños de alta especificidad
lingüística, la palanca correcta es ampliar la capacidad del lado
visual sobre el propio dataset, no incorporar imágenes externas
estilísticamente discordantes.

\subsection{Calibración por \emph{temperature scaling}}
\label{sec:anexoi-calibracion}

Los logits \emph{zero-shot} de la rama contrastiva están severamente
sobreconfiados, propiedad esperada en arquitecturas CLIP-style con
InfoNCE. La calibración con \emph{temperature scaling} de un parámetro
escalar~\cite{guo2017} sobre los $107$ pares de validación reduce el
ECE sobre los $101$ pares de test
(tabla~\ref{tab:spanderm-calibration},
figura~\ref{fig:spanderm-calibration}).

\begin{table}[H]
\centering
\footnotesize
\setlength{\tabcolsep}{4pt}
\renewcommand{\arraystretch}{0.95}
\caption{Calibración post-hoc por \emph{temperature scaling} sobre
la iteración contrastiva v4 de Dermapixel R0. Ajuste del escalar
$\tau$ por optimización LBFGS
sobre $107$ pares de validación; ECE medido sobre $101$ pares de
test. La calibración sobre v5 no converge (la optimización LBFGS
diverge sobre el espacio embebido \emph{image-LoRA} expandido y
$\tau \to 0$), por lo que queda registrada como limitación
metodológica abierta.}
\label{tab:spanderm-calibration}
\begin{tabular}{lrrr}
\hline
Nivel & ECE antes  & ECE después & $\tau$       \\
\hline
L3 (v4) & $0{,}7262$ & $0{,}0213$  & $0{,}4165$ \\
\hline
\end{tabular}
\end{table}

\begin{figure}[H]
\centering
\includegraphics[width=0.88\textwidth]{fig_spanderm_calibration.png}
\caption{Diagrama de fiabilidad sobre L3 antes y después de
\emph{temperature scaling} ($\tau = 0{,}4165$) sobre la iteración
contrastiva v4 de Dermapixel R0. ECE desciende de $0{,}7262$ a
$0{,}0213$ ($-97\%$). La
calibración escalar de un parámetro restaura la interpretabilidad
probabilística de los logits sin alterar el \emph{ranking} entre
clases.}
\label{fig:spanderm-calibration}
\end{figure}

La reducción del ECE es del $-97\%$ sobre L3 con un único parámetro
escalar. Convergente con la sensibilidad de DermLIP
(sección~\ref{sec:zs-results}): sin calibración post-hoc, las
probabilidades de un sistema CLIP-style \emph{zero-shot} castellano no
son interpretables como confianza clínica, por lo que el despliegue
debe incorporar este escalado como capa final ajustado sobre el
archivo de destino. La extensión a v5 no converge ($\tau \to 0$ sobre
el espacio \emph{image-LoRA} expandido), limitación que no afecta al
resultado principal de v4.

\subsection{Diagnóstico cuantitativo del techo sobre L2}
\label{sec:anexoi-l2-ceiling}

La tabla~\ref{tab:spanderm-l2-ceiling} sitúa la iteración contrastiva
v2 en \emph{zero-shot} L2 frente a las dos referencias internas del
capítulo de resultados (sección~\ref{sec:dermapixel-results}
y~\ref{sec:dermapixel-spanderm-v0}).

\begin{table}[H]
\centering
\footnotesize
\setlength{\tabcolsep}{4pt}
\renewcommand{\arraystretch}{0.95}
\caption{Rama contrastiva de Dermapixel R0 \emph{zero-shot} L2
frente a las dos
referencias internas del capítulo sobre el mismo nivel ontológico.
Las tres referencias evalúan sobre el conjunto de test de
DermapixelAI~1.0; la cifra de la iteración v2 corresponde al
test de $101$ imágenes del \emph{split} per-case del experimento
contrastivo y las dos restantes al test fijo
$N_{\mathrm{test}} = 36$ de la sección~\ref{sec:dermapixel-results}.
La comparación es indicativa por la diferencia muestral.}
\label{tab:spanderm-l2-ceiling}
\adjustbox{max width=\textwidth}{%
\begin{tabular}{llrr}
\hline
Sistema                                                 & Régimen               & Acc@1 L2          & Observaciones \\
\hline
LP PanDerm Large (referencia, tabla~\ref{tab:dermapixel-lp}) & Supervisado L2 cerrado & $0{,}250$        & $38$ clases, $N_{\mathrm{test}} = 36$ \\
Dermapixel R0 cabeza L2 (tabla~\ref{tab:dermapixel-spanderm-l2}) & Supervisado L2 cerrado & $0{,}250$        & $38$ clases, $N_{\mathrm{test}} = 36$ \\
\textbf{Dermapixel R0 contrastivo v2}                    & \emph{Zero-shot} open  & $\mathbf{0{,}133}$ & $23$ clases efectivas, $N_{\mathrm{test}} = 101$ \\
\hline
\end{tabular}}
\end{table}

La jerarquía entre regímenes es la esperada: la cabeza supervisada
cerrada con LoRA (Dermapixel R0, cabeza L2) alcanza Acc@1 $0{,}250$
sobre L2, frente a $0{,}133$ de la rama contrastiva \emph{zero-shot
open-class} (v2) sobre las $23$ clases efectivas del test per-case. La
diferencia operativa es la capacidad de la rama contrastiva de aceptar
consultas textuales libres en castellano sin fijar el vocabulario. Esa
propiedad motivó un futuro buscador textual del prototipo, pero no se
despliega ---su recuperación no generaliza al uso real
(sección~\ref{sec:anexoi-sintesis})---; la búsqueda de casos similares en
producción la resuelve M4-bis (imagen$\to$imagen, PanDerm Large $+$ FAISS;
Anexo~\ref{cap:anexoh}).

\subsection{Síntesis de las cinco iteraciones}
\label{sec:anexoi-sintesis}

La tabla~\ref{tab:spanderm_iteraciones} sintetiza la trayectoria
completa de la rama contrastiva, articulando por iteración la
configuración LoRA, el régimen de inferencia, las dos cifras de
referencia (L3 \emph{accuracy}@1 \emph{zero-shot} y lift de
recuperación $i2t$ R@$5$) y la lectura asociada.

\begin{table}[H]
\centering
\footnotesize
\setlength{\tabcolsep}{3pt}
\renewcommand{\arraystretch}{1.05}
\caption{Comparativa de las cinco iteraciones del entrenamiento
contrastivo de Dermapixel R0 sobre el dataset DermapixelAI~1.0. Régimen evaluativo
\emph{zero-shot} L3 sobre $101$ imágenes de test ($54$ clases L3
representadas), \emph{H6 hybrid} v$30$\_r$70$ (combinación
$30\%$ prototipos visuales $+$ $70\%$ texto rico) aplicado de
forma homogénea a las cinco iteraciones. v4 queda a $0{,}0025$ del
objetivo fijado a priori de $0{,}25$ acc@1 L3 (dentro del ruido de
muestreo) mediante expansión arquitectónica
(\emph{image-LoRA} $r=32$ sobre $12$ bloques $+$ \emph{ensemble}
híbrido texto-visual). v5 confirma negativamente que la escala
externa genérica con \emph{captions} sintéticas no supera v4. Para
v1--v3 los prototipos visuales derivan del espacio embebido del
\emph{checkpoint} sin LoRA visual; v4--v5 emplean prototipos
sobre el espacio embebido con \emph{image-LoRA} expandido. La
comparativa permanece interpretable porque los prototipos derivan
del propio espacio embebido del \emph{checkpoint} evaluado.
Provenance: commit \texttt{d0190c3} del repositorio
\texttt{spanderm-clip}. Las cifras Linear Probing BAcc L2 absolutas
para v4 ($0{,}1860$) y v5 ($0{,}2871$) proceden de los ficheros
JSON \texttt{eval\_v5/linear\_probe\_v4ref.json} y
\texttt{eval\_v5/linear\_probe\_v5.json}; el delta relativo
$+54{,}4\%$ queda confirmado.}
\label{tab:spanderm_iteraciones}
\adjustbox{max width=\textwidth}{%
\begin{tabular}{lllrrll}
\hline
Versión & LoRA texto & LoRA visual & acc@1 L3 & Lift i2t R@$5$ & LP BAcc L2 & Notas \\
\hline
v1 & ---                       & ---                          & $0{,}1782$          & ---                       & ---                                 & Baseline e5 \emph{full FT} \\
v2 & $r=16$, $\alpha=32$       & ---                          & $0{,}1980$          & $\mathbf{14{,}0\times}$   & ---                                 & Pico retrieval rich \\
v3 & $r=16$, $\alpha=32$       & $r=16$, últ.~$4/24$           & $0{,}1980$          & ---                       & ---                                 & Image-LoRA infradotado \\
v4 & $r=32$, $\alpha=64$       & $r=32$, últ.~$12/24$          & $\mathbf{0{,}2475}$ & $12{,}4\times$            & $0{,}1860$   & A $0{,}0025$ del objetivo \\
v5 & $r=32$, $\alpha=64$       & $r=32$, últ.~$12/24$          & $0{,}2079$          & $10{,}8\times$            & $0{,}2871$   & FLAT: escala externa ($+54{,}4\%$ relativo sobre v4) \\
\hline
\end{tabular}}

\smallskip
\footnotesize
\textit{Nota histórica.} Los valores \emph{rich-only baseline}
($\sim 30$--$80$ palabras/clase) fueron $0{,}1287$ (v1), $0{,}1386$
(v2) y $0{,}1089$ (v3); la transición al régimen \emph{H6 hybrid}
homogéneo de esta tabla permite la comparativa cross-\emph{checkpoint}
(detalle en Anexo~\ref{cap:anexog}).
\normalsize
\end{table}

\textbf{Importante:} las cifras de \emph{Lift i2t} R@$5$ de la tabla
(con pico $14{,}0\times$ en v2, equivalente a R@$5 = 0{,}654$)
corresponden al régimen de \emph{validación} con descripciones ricas
(\texttt{llm\_case\_summary}) y caracterizan un escenario favorable, no el
rendimiento del buscador en uso real. Sobre el test \emph{held-out} y con
consulta corta, la recuperación cae a niveles de azar (i2t R@$5$ del orden
de $0{,}02$), por lo que estas cifras deben leerse como evidencia de la
\emph{viabilidad técnica} del espacio dual texto-imagen, no como capacidad
de recuperación desplegable.

El experimento articula tres ejes. La capacidad arquitectónica
(v3$\to$v4) es la palanca efectiva: el cuello de botella reside en el
\emph{image encoder}, no en la formulación textual. La escala externa
(v5) queda \emph{flat-to-negative}, convergente con la sensibilidad de
DermLIP al tokenizador (sección~\ref{sec:zs-results}). El eje de
supervisión conceptual ---anotaciones del \emph{Seven-Point
Checklist}~\cite{kawahara2019} o del Grupo~A (Anexo~\ref{cap:anexof})---
queda abierto como línea de aprendizaje activo dirigido
(sección~\ref{sec:linea-active-learning}). El experimento valida así la
utilidad de PanDerm Large para tareas castellanas más allá del LP/FT
supervisado (v4 a $0{,}0025$ del umbral objetivo L3 $0{,}25$) y
establece una metodología reproducible para evaluaciones CLIP-style
sobre archivos privados de pequeña escala.
